% --- Back cover ---
% Large PERSONIX wordmark, no logo image. Cream paper inherited globally.
% Absolute (tikz overlay) positioning so the footer sticks to the bottom.
\clearpage
\thispagestyle{empty}
\begin{tikzpicture}[remember picture, overlay]
  % --- Wordmark ---
  \node[anchor=north, font=\sffamily\bfseries\fontsize{48}{52}\selectfont]
    at ([yshift=-26mm]current page.north) {PERSONIX};
  \node[anchor=north, font=\rmfamily\itshape\large]
    at ([yshift=-44mm]current page.north) {Schimbare fără compromis};
  \node[anchor=north, font=\rmfamily\small,
        text width=\paperwidth-50mm, align=center]
    at ([yshift=-54mm]current page.north)
    {O rețea de reputație descentralizată, necenzurabilă și incoruptibilă};
  % --- Body ---
  \node[anchor=north, text width=\paperwidth-50mm, align=justify, font=\small]
    at ([yshift=-72mm]current page.north)
    {%
      Statul funcționa când comunicarea era lentă, deciziile erau locale, iar
      autoritatea trăia în același oraș pe care îl guverna. Rețelele de azi
      inversează toate cele trei presupuneri --- iar instituțiile clădite pe
      ele nu se mai potrivesc lumii pe care o guvernează. Această carte descrie
      un instrument care se potrivește: o rețea de reputație descentralizată în
      care identitatea îți rămâne ție, adevărul este auditabil public, iar
      mituirea este o sinucidere reputațională.

      \smallskip
      Nu un manifest. Nu o prognoză. Un plan de lucru pentru cum poate fi
      construit succesorul statului --- de bunăvoie, câte o declarație pe rând.
    };
  % --- Footer (anchored to the bottom of the page) ---
  \node[anchor=south, font=\sffamily\footnotesize,
        text width=\paperwidth-50mm, align=center]
    at ([yshift=12mm]current page.south)
    {%
      Pavel Kudrna\quad\textbullet\quad Praga, 2026\par
      \href{https://personix.org}{personix.org}\par
      Licențiat sub Creative Commons Attribution-ShareAlike 4.0 International
      (CC BY-SA 4.0).
    };
\end{tikzpicture}%
\mbox{}%
\clearpage
